\chapter{Running the Code}
\label{app:reproducibility}

This appendix gives the practical steps needed to reproduce the reported
experiments. It is a compact execution guide for setting up the repository,
running the solver, auditing output files, and regenerating the reported tables
and figures.

\section{System and Software Requirements}

The reported experiments were run in a CPU-only local environment on macOS with
an Apple M3 processor, 8 CPU cores, and 16 GB memory. No GPU, HPC cluster, or
commercial optimization solver was required for the reported runs. The solver
requires Python 3.11 or newer and the packages listed in
\texttt{requirements.txt}. Java is required only for the external MINOA
feasibility and objective audit.

\begin{lstlisting}[basicstyle=\ttfamily\scriptsize,breaklines=true,aboveskip=4pt,belowskip=4pt]
python3 -m venv .venv
source .venv/bin/activate
python -m pip install --upgrade pip
python -m pip install -r requirements.txt
java -version
\end{lstlisting}

\section{Repository Structure}

The repository is organized as a code-first project. The most important tracked
folders and files are:

\begin{lstlisting}[basicstyle=\ttfamily\scriptsize,breaklines=true,aboveskip=4pt,belowskip=4pt]
project-root/
  .github/workflows/                  CI checks for headline and all-instance runs
  data/raw/minoa/senior/              Senior benchmark input files
  scripts/                            command-line runners and figure scripts
  scripts/minoa_lib/                  implementation modules
  tests/                              rule-focused tests
  tools/minoa/desktopValidator/...jar external MINOA checker
  results/final_validated_results.json canonical thesis result table
  requirements.txt                    Python dependencies
\end{lstlisting}

Generated working files are written outside the tracked source structure, mainly
under \texttt{outputs/minoa/} and \texttt{data/processed/minoa/}. The raw Senior
input files are not modified by the solver.

The public repository and the continuous-integration runs are available through
the \href{https://github.com/mirza019/minoa-mixed-fleet-solver}{project
repository}, the \href{https://github.com/mirza019/minoa-mixed-fleet-solver/actions}{GitHub
Actions overview}, the
\href{https://github.com/mirza019/minoa-mixed-fleet-solver/actions/workflows/minoa-headline-instances.yml}{headline
workflow}, the
\href{https://github.com/mirza019/minoa-mixed-fleet-solver/actions/workflows/minoa-all-instances.yml}{all-instance
workflow}, and the
\href{https://github.com/mirza019/minoa-mixed-fleet-solver/actions/workflows/results-check.yml}{result
consistency workflow}.

\section{Main Experiment Commands}

The main entry point is \texttt{scripts/run\_experiment.py}. The final thesis
method is selected with \texttt{--algorithm multistart}. The following commands
cover the main reproduction workflow:

\begin{lstlisting}[basicstyle=\ttfamily\scriptsize,breaklines=true,aboveskip=4pt,belowskip=4pt]
# Headline instances: Small, Medium, Large
.venv/bin/python scripts/run_experiment.py --algorithm multistart --scope sml

# One headline instance
.venv/bin/python scripts/run_experiment.py --algorithm multistart --scope sml --only Small

# Final all-instance no-regression archive
.venv/bin/python scripts/run_experiment.py --algorithm multistart --scope all

# Optional fresh direct all-instance audit before printing the final archive
.venv/bin/python scripts/run_experiment.py --algorithm multistart --scope all --fresh-audit
\end{lstlisting}

The all-instance archive command prints the canonical final thesis table with
total validated cost \(10000.48\) and \(126\) used vehicle blocks. The
\texttt{--fresh-audit} option reruns the construction pipeline first and then
prints the retained no-regression archive.

The comparison algorithms use the same command style:

\begin{lstlisting}[basicstyle=\ttfamily\scriptsize,breaklines=true,aboveskip=4pt,belowskip=4pt]
.venv/bin/python scripts/run_experiment.py --algorithm greedy --scope sml
.venv/bin/python scripts/run_experiment.py --algorithm pathcover --scope sml
.venv/bin/python scripts/run_experiment.py --algorithm weighted --scope sml
.venv/bin/python scripts/run_experiment.py --algorithm multistart --scope sml
\end{lstlisting}

\section{Output Files and External Audit}

Headline outputs are written to \texttt{outputs/minoa/sml\_<algorithm>/}. The
final retained archive is stored under \texttt{outputs/minoa/final\_archive/}.
The compact machine-readable result file is
\texttt{results/final\_validated\_results.json}.

A generated input-output pair can be audited manually with the external Java
checker:

\begin{lstlisting}[basicstyle=\ttfamily\scriptsize,breaklines=true,aboveskip=4pt,belowskip=4pt]
java -jar tools/minoa/desktopValidator/desktopValidator/desktopValidator.jar \
  data/raw/minoa/senior/Small_Input_S.json \
  outputs/minoa/sml_multistart/Small_Output_multistart.json
\end{lstlisting}

The audit checks feasibility and reports the objective value for the complete
input-output pair. It is not an optimality certificate. Existing output files
can be summarized in a table with:

\begin{lstlisting}[basicstyle=\ttfamily\scriptsize,breaklines=true,aboveskip=4pt,belowskip=4pt]
.venv/bin/python scripts/minoa_report.py INPUT_JSON:OUTPUT_JSON
\end{lstlisting}

\section{Tables, Figures, and Lower Bounds}

The final compact thesis result table is generated from
\texttt{results/final\_validated\_results.json}:

\begin{lstlisting}[basicstyle=\ttfamily\scriptsize,breaklines=true,aboveskip=4pt,belowskip=4pt]
.venv/bin/python scripts/print_final_results_table.py
\end{lstlisting}

The numerical result graphs are regenerated with:

\begin{lstlisting}[basicstyle=\ttfamily\scriptsize,breaklines=true,aboveskip=4pt,belowskip=4pt]
.venv/bin/python scripts/minoa_make_constraint_figures.py
\end{lstlisting}

Lower-bound reports are separate from feasible schedule construction:

\begin{lstlisting}[basicstyle=\ttfamily\scriptsize,breaklines=true,aboveskip=4pt,belowskip=4pt]
.venv/bin/python scripts/run_lower_bounds.py \
  --scope all \
  --input-dir data/processed/minoa/final_adaptive_bounded \
  --archive-csv outputs/minoa/final_archive/final_results.csv \
  --time-limit 180 \
  --output-csv results/lower_bounds/all_instances_lower_bounds.csv

.venv/bin/python scripts/generate_lower_bound_figures.py \
  --csv results/lower_bounds/all_instances_lower_bounds.csv \
  --out-dir outputs/minoa/thesis_figures
\end{lstlisting}

Only certified global lower bounds are used as global bounds. Selected-timetable
bounds are diagnostic values for a fixed timetable and are not global optimality
certificates.

\section{Quick-Start Recipe}

The following command sequence is sufficient for a quick reproduction run:

\begin{lstlisting}[basicstyle=\ttfamily\scriptsize,breaklines=true,aboveskip=4pt,belowskip=4pt]
python3 -m venv .venv
source .venv/bin/activate
python -m pip install -r requirements.txt

.venv/bin/python scripts/run_experiment.py --algorithm multistart --scope sml
.venv/bin/python scripts/run_experiment.py --algorithm multistart --scope all
.venv/bin/python scripts/print_final_results_table.py
\end{lstlisting}

For a manual audit of a specific output file, use the Java checker command from
the previous section.

\section{Troubleshooting}

\begin{table}[H]
\centering
\scriptsize
\caption{Common execution issues and practical fixes.}
\label{tab:appendix-troubleshooting}
\begin{tabularx}{\textwidth}{p{3.8cm}X}
\toprule
Issue & Practical fix \\
\midrule
Python package missing &
Activate \texttt{.venv} and run \texttt{python -m pip install -r requirements.txt}. \\
Java not found &
Run \texttt{java -version}. If it fails, install a Java Runtime Environment or JDK. \\
Input file not found &
Run commands from the project root and check \texttt{data/raw/minoa/senior/}. \\
Output folder missing &
Use the experiment wrappers; they create output folders automatically. \\
Checker path changed &
Pass the checker location with \texttt{--validator}. \\
Large run is slow &
Use \texttt{--scope sml} for a quick check. Use \texttt{--quick} only for diagnostics,
not for final reported results. \\
\bottomrule
\end{tabularx}
\end{table}
